% 05-Results-prose.tex -- observation paragraphs to sit alongside the
% exhibits in 05-Results.tex. Drafted for the professor to write from:
% each paragraph states what an exhibit shows and nothing beyond it.
%
% AAAI-safe: no packages, no forbidden commands, pure ASCII.
%
% Suggested placement is marked before each block. The exhibits are
% unchanged; this file adds only prose.

% ---------------------------------------------------------------- opening
% Place immediately after \section{Results}, before Table 1.

Across the three trade years the island-champion ensemble returns
$+27.5\%$ at the registered 40-island width, against $+11.3\%$ for an
online LSTM, $+11.7\%$ for an online GRU and $+14.8\%$ for an LSTM
retrained monthly (Table~\ref{tab:window}). Paired within each (panel,
seed) cell the margins are $+16.2$, $+15.8$ and $+12.7$ points, at
$t=7.3$, $6.6$ and $6.7$ over 40 cells (Table~\ref{tab:paired}). The
ordering is stable year by year: the ensemble leads in each of 2022,
2023 and 2024 rather than earning its margin in a single favourable year.

% ------------------------------------------------- the central observation
% Place after Table 1 / Figure 1, as the paper's main claim.

The more informative comparison is inside the system rather than against
the baselines. The single global-best genome --- the prediction rule of
prior ONE-NAS work --- returns $+4.5\%$ at 40 islands, $+4.7\%$ at 20 and
$+1.0\%$ at 60 (Table~\ref{tab:window}). Those figures sit below every
neural baseline and below buy and hold. The same runs, read as an
ensemble of island champions rather than as a single champion, return
$+27.5$, $+27.4$ and $+30.1$. The search therefore does not pay by
finding one good architecture; it pays by producing a population.

Table~\ref{tab:m1} makes the comparison within-run and paired, since both
rules are scoring-time readings of the same evolved population: the
ensemble adds $+23.0$ points at 40 islands ($t=18.2$), $+22.6$ at 20
($t=14.9$) and $+29.1$ at 60 ($t=18.3$), with Sharpe gains of $+0.62$ to
$+0.84$ on the same pairing. Nothing about the search changes between the
two columns --- only which genomes are read.

% ------------------------------------------------------------- mechanism
% Place after the diversity table.

The size of that gain is not an empirical surprise once member diversity
is measured. Table~\ref{tab:diversity} reports, for each width, the mean
rank IC of an individual island champion, the mean pairwise Spearman
correlation $\rho$ between member cross-sections, and the resulting
ensemble IC. Members are individually weak (rank IC $+0.0065$ to
$+0.0068$) and substantially decorrelated ($\rho$ between $0.150$ and
$0.166$). The equicorrelated averaging prediction
$m\sqrt{N/(1+(N-1)\rho)}$ then anticipates the measured ensemble IC to
within 3\% at every width. The ensemble is not recovering a latent
strong model; it is averaging many weak, weakly-correlated ones, and the
observed gain is the size that diversity implies.

% ------------------------------------------------------------- invariance
% Place after the HP-invariance table.

The effect does not depend on the configuration it was measured in.
Table~\ref{tab:invariance} recomputes the ensemble-minus-champion
difference inside all 16 trained hyperparameter configurations of the
screen --- different selection metrics, target horizons, population
ratios, repopulation frequencies, replay settings and training budgets.
The difference is positive in 16 of 16, with a mean of $+8.3$ points.
It also holds at every island width tested and on the tuning span as well
as the evaluation span (Table~\ref{tab:m1}).

% ------------------------------------------------------------------ width
% Place after the island-count figure.

Figure~\ref{fig:islands} separates two things that island count does.
Return and Sharpe rise steeply from 10 to about 20 islands ($+19.9$ to
$+27.3$ net, Sharpe $0.88$ to $1.11$) and then flatten: across 20 to 60
the pooled net sits between $+26$ and $+30$ and the widths are not
separable at ten seeds. Seed dispersion narrows over the same range, the
standard deviation of per-seed Sharpe falling from $0.227$ at 10 islands
to $0.11$--$0.13$ above 30 and the worst seed rising from $0.53$ to
$0.91$, though not monotonically width by width. Width therefore buys
consistency at least as much as expected return, which is the relevant
property for a deployed system, where a single run is held rather than an
average over runs.

Width was chosen on the tuning span, not on this curve.
Table~\ref{tab:tunewidth} gives the same sweep on 2016--2019: rank IC
rises from 8 to about 20 islands and is flat above it, with the 40-, 50-
and 60-island cells spanning $0.0168$--$0.0170$. Measured on the
evaluation span 50 and 60 islands score above 40; measured on the tuning
span, where a configuration choice can legitimately be made, they are
indistinguishable from it (largest $|t|=0.23$). The apparent width
advantage does not survive protocol-symmetric selection and is not
adopted.

% ------------------------------------------------------------ risk, costs
% Place after the alphas and cost tables.

The returns are not explained by factor exposure. Regressed on the
Fama--French three factors plus momentum with Newey--West standard errors
(Table~\ref{tab:alphas}), all three ONE-NAS widths carry alphas
distinguishable from zero ($+10.0$ to $+11.2\%$/yr, $t=2.56$ to $2.77$)
at market betas near $0.11$--$0.12$. Among the baselines only the
monthly-retrained LSTM reaches $|t|=2$, at $+6.0\%$/yr and $t=2.01$;
buy and hold earns no alpha at all ($-2.3\%$/yr, $t=-0.72$) at a market
beta of $0.80$.

Nor do the results depend on an optimistic cost assumption.
Table~\ref{tab:costs} scales the realised per-name transaction cost while
preserving its cross-sectional shape. The ensemble remains profitable to
roughly $10$--$12\times$ realistic costs, against $4.5$--$5.4\times$ for
the baselines, because it also turns over least.

% ------------------------------------------------------------ limitations
% Place at the end of Results, or fold into the Discussion.

Three limits are worth stating with the results.

The reported window is a favourable one for the comparison against buy
and hold, and the full scored span is not. On 2022--2024 an equal-weight
hold of the same universe returns $+7.9\%$ under the prior study's
convention, against $+27.5\%$ for the ensemble. Scored instead over the
whole registered span, 2020--2024, the same benchmark returns $+44.3\%$
against the ensemble's $+45.6\%$ --- a tie on raw return, because
2020--2021 was a strong market that a long-only hold captures and a
dollar-neutral book does not. The defensible claim is therefore
risk-adjusted rather than a return margin: the ensemble matches or beats
the benchmark while running at a market beta near $0.11$ rather than
$0.80$, and it is the only construction here with an alpha
distinguishable from zero.

Second, the advantage is regime-dependent. In 2020, the crisis year
excluded from the reported window, a plain online LSTM adapts faster than
the evolved population and outperforms it; the ensemble's edge is largest
in ordinary markets.

Third, panel membership is fixed over the sample, so the universe is
conditioned on names that survived to the end of it. This applies equally
to every arm compared here, but it does mean the absolute return levels
should not be read as achievable ex ante.
